% M07, 11.09.2026: F1 als Leitbeispiel früh eingeführt; Zusammenhang
% von Designanforderungen, Architektur und MAF-Realisierung erhalten.
% F1: synthetisches meeting-4-ux-block-l.txt; ausgeführter Folgefall
% 20260818_084712_765eea, bestehender Core, Apply-Bericht zu Issues 44/45.

\chapter{Lösungskonzeption und Architektur des Artefakts}
\label{chap:loesungsarchitektur}

\section{Zielbild, Anforderungen und Systemgrenze}
\label{sec:zielbild-anforderungen-systemgrenze}

Das Artefakt überführt wiederkehrende Projektkommunikation in
quellengebundene Anforderungen, Architekturinformationen und
Arbeitspakete. Neue Inhalte werden im laufenden Betrieb gegen den
vorhandenen Projektbestand eingeordnet. Die Anforderungen A1 bis A10
aus Abschnitt~\ref{sec:konsolidierte-designanforderungen} begründen
die Aufgabenverteilung: Agentische Komponenten interpretieren
Inhalte und erzeugen Vorschläge; explizite Workflows steuern deren
Prüfung und Weitergabe. Die Fortschreibung des Projektbestands
unterliegt deterministischen Integritätsprüfungen und menschlicher
Autorisierung. Der Evidence Ledger strukturiert die
transkriptbasierte Evidenz, der Core führt den maßgeblichen
fachlichen Projektzustand über einzelne Läufe hinweg. GitHub und
Dokumente bilden daraus abgeleitete Sichten. Die Zuordnung der
Anforderungen zu den sechs Architektursäulen zeigt
Tabelle~\ref{t:saeulen-uebergabe}.

Als wiederkehrendes Beispiel dient die \emph{Besuchsankündigung}
für die in Abschnitt~\ref{sec:ausgangslage-entwicklungsbefunde}
eingeführte App: Angehörige sollen Besuche mit Datum und Uhrzeit
ankündigen können; Pflegende benötigen eine Übersicht. Dieser
Wunsch stammt aus einem gezielt synthetisch erstellten Folgegespräch,
das im dokumentierten Fall F1 tatsächlich verarbeitet wurde. Der
Fall beginnt mit einem vorhandenen Projektbestand und führt bis zu
real erzeugten GitHub-Issues. Kurze Wiederaufnahmen erläutern die
jeweilige Verarbeitung; Abschnitt~\ref{sec:endtoend-pfade} führt die
Stationen zusammen und unterscheidet diesen Betriebspfad von der
initialen Bestandsherstellung.

Der betrachtete SDLC-Ausschnitt endet bei strukturierter
Arbeitsplanung und Außenprojektion; Codegenerierung und
Implementierung der geplanten App liegen außerhalb.
Das Forschungsartefakt adressiert einen Einzelbediener-Kontext;
Bedienqualität und Produktivität werden nicht bewertet. Die
Außenkopplung ist für GitHub realisiert. Dieses Kapitel erklärt
die frameworkneutrale Architektur; Kapitel~\ref{chap:maf-realisierung}
zeigt ihre konkrete MAF-Realisierung. Kapitel~\ref{chap:evaluation}
untersucht ausgewählte Eigenschaften und Grenzen anhand der dort
ausgewiesenen Quellen, Fälle und Systemstände.
